\begin{longtblr}
[
 label={tab:cierre_vacios_implicaciones},
 caption={Cierre analítico: vacíos recurrentes, impacto metodológico e implicaciones operativas. Elaboración propia.}]{
  colspec={X[2.2,l,m] X[0.8,c,m] X[0.8,c,m] X[0.9,c,m] X[2.6,l,m] X[2.8,l,m]},
  row{1}={font=\bfseries}
}
Vacío/Límite recurrente & k & n & \% & Impacto en síntesis & Implicación para la propuesta \\
Partición de datos no reportada \texttt{NR} & 10 & 19 & 52.6 & Introduce incertidumbre en la comparación de desempeño por variación no controlada del esquema de validación. & Definir y justificar un esquema fijo de evaluación \texttt{train/val/test} o validación cruzada), con trazabilidad completa de su aplicación. \\
Defuzzificación no reportada & 8 & 19 & 42.1 & Dificulta la reproducibilidad del núcleo inferencial y la atribución técnica del desempeño al diseño difuso. & Especificar método de defuzzificación y parametrización asociada como componente obligatorio del pipeline. \\
Comparador ausente \texttt{NR} & 4 & 19 & 21.1 & Reduce la validez comparativa de afirmaciones de mejora frente a enfoques alternativos. & Incluir un benchmark mínimo y un criterio homogéneo de comparación cuantitativa en la experimentación. \\
Mejora no reportada \texttt{NR} & 4 & 19 & 21.1 & Limita la interpretación del aporte incremental al no explicitar la magnitud de ganancia. & Reportar mejora absoluta y relativa sobre baseline, bajo métricas y condiciones equivalentes. \\
Heterogeneidad de métricas (sin núcleo único) & 11 & 19 & 57.9 & Fragmenta la lectura transversal e incrementa heterogeneidad de reporte entre estudios. & Establecer un conjunto base de métricas \texttt{MSE/RMSE/MAE} y, cuando aplique, \texttt{Accuracy} de reporte obligatorio. \\
\end{longtblr}
